\documentclass[11pt]{article}
\usepackage[margin=1in]{geometry}
\usepackage{hyperref}
\usepackage{listings}
\usepackage{xcolor}
\usepackage{enumitem}

\lstset{
    basicstyle=\ttfamily\small,
    backgroundcolor=\color{gray!10},
    frame=single,
    framerule=0.5pt,
    breaklines=true,
    showstringspaces=false,
    language=bash,
    morekeywords={cmake,ninja,git,cd,mkdir,set}
}

\title{Building femm-qt on Windows}
\author{D.\ Meeker}
\date{\today}

\begin{document}
\maketitle

\section{Overview}

femm-qt is a cross-platform port of FEMM (Finite Element Method Magnetics)
using Qt~6 for the GUI and CMake/Ninja for the build system.  This document
describes how to build it from source on Windows using the Qt Online
Installer and MinGW (no Visual Studio required).

The build produces six executables:
\begin{itemize}[nosep]
  \item \texttt{femm\_gui.exe} --- the main GUI application
  \item \texttt{femm\_headless.exe} --- headless IPC server (no GUI)
  \item \texttt{fkn.exe} --- magnetics solver
  \item \texttt{belasolv.exe} --- electrostatics solver
  \item \texttt{hsolv.exe} --- heat flow solver
  \item \texttt{csolv.exe} --- current flow solver
\end{itemize}

\section{Prerequisites}

\subsection{Qt 6 with MinGW}

Download and run the Qt Online Installer from
\url{https://www.qt.io/download-qt-installer-oss}.  You will need a
free Qt account.

In the installer, select:
\begin{itemize}[nosep]
  \item A recent Qt version (e.g.\ Qt 6.10.3)
  \item Under that version, check \texttt{MinGW 64-bit}
  \item Under \textbf{Developer and Designer Tools}, check:
    \begin{itemize}[nosep]
      \item \texttt{MinGW 13.1.0 64-bit} (or the version matching your Qt)
      \item \texttt{CMake}
      \item \texttt{Ninja}
    \end{itemize}
\end{itemize}

The default install location is \texttt{C:\textbackslash Qt}.  The paths
below assume this default; adjust if you installed elsewhere.

After installation, verify the tools are accessible.  Open a command
prompt and check:
\begin{lstlisting}
C:\Qt\Tools\mingw1310_64\bin\g++ --version
C:\Qt\Tools\CMake_64\bin\cmake --version
C:\Qt\Tools\Ninja\ninja --version
\end{lstlisting}

\subsection{Git (Optional)}

If you are cloning from a repository rather than using a source zip,
install Git for Windows from \url{https://git-scm.com/download/win}.

\section{Getting the Source}

If you have a source zip, extract it:
\begin{lstlisting}
unzip femm-qt-src.zip
cd femm-qt
\end{lstlisting}

Or clone from a repository:
\begin{lstlisting}
git clone <repository-url> femm-qt
cd femm-qt
\end{lstlisting}

\section{Building}

From the \texttt{femm-qt} directory, run the build script:
\begin{lstlisting}
win_build.bat
\end{lstlisting}

This auto-detects Qt, MinGW, CMake and Ninja under \texttt{C:\textbackslash Qt},
configures the project with CMake, and builds it with Ninja.  All
executables are placed in \texttt{build\textbackslash bin\textbackslash}.

To wipe the build directory and rebuild from scratch, pass
\texttt{clean}:
\begin{lstlisting}
win_build.bat clean
\end{lstlisting}

If the script cannot find your Qt or MinGW installation --- for
example, if Qt is installed somewhere other than \texttt{C:\textbackslash Qt}
--- see \textit{Manual configuration} under Troubleshooting.

\section{Running}

Launch the GUI:
\begin{lstlisting}
bin\femm_gui.exe
\end{lstlisting}

Or the headless server:
\begin{lstlisting}
bin\femm_headless.exe
\end{lstlisting}

\subsection{Setting FEMM\_DIR}

The language bindings (OctaveFEMM, MathFEMM, pyFEMM, JuliaFEMM) locate
the FEMM executables via the \texttt{FEMM\_DIR} environment variable.
Set it to the \texttt{bin} directory:
\begin{lstlisting}
set FEMM_DIR=C:\path\to\femm-qt\build\bin
\end{lstlisting}
To make this persistent, add it via System Properties $\rightarrow$
Environment Variables, or in PowerShell:
\begin{lstlisting}
[Environment]::SetEnvironmentVariable("FEMM_DIR",
    "C:\path\to\femm-qt\build\bin", "User")
\end{lstlisting}

\section{Using with MATLAB or Octave}

\subsection{MATLAB}

Set \texttt{FEMM\_DIR} before launching MATLAB, or set it within MATLAB:
\begin{lstlisting}
setenv('FEMM_DIR', 'C:\path\to\femm-qt\build\bin')
\end{lstlisting}
Add the \texttt{octavefemm} directory to the MATLAB path:
\begin{lstlisting}
addpath('C:\path\to\femm-qt\octavefemm')
\end{lstlisting}
Then call \texttt{openfemm} to launch femm-qt.

\subsection{Octave}

Octave on Windows requires a modern JDK for socket communication
(Adoptium Temurin 25 or later is recommended; MATLAB's bundled Java~8
may segfault with Octave).  Set the \texttt{JAVA\_HOME} environment
variable to point at the JDK, then launch Octave.

\begin{lstlisting}
set JAVA_HOME=C:\Program Files\Eclipse Adoptium\jdk-25
\end{lstlisting}

Add the \texttt{octavefemm} directory to the Octave path and call
\texttt{openfemm}.

\section{Troubleshooting}

\subsection{Manual configuration}

\texttt{win\_build.bat} auto-detects the toolchain under
\texttt{C:\textbackslash Qt}.  If it picks the wrong Qt version, or your
tools are installed elsewhere, set the \texttt{QT\_DIR} and
\texttt{MINGW\_DIR} variables near the top of the script to the correct
paths and run it again.

To configure entirely by hand instead, create the build directory and
invoke CMake and Ninja directly:
\begin{lstlisting}
mkdir build
cd build
cmake .. -G Ninja -DCMAKE_BUILD_TYPE=Release ^
    -DCMAKE_PREFIX_PATH="C:/Qt/6.10.3/mingw_64" ^
    -DCMAKE_C_COMPILER="C:/Qt/Tools/mingw1310_64/bin/gcc.exe" ^
    -DCMAKE_CXX_COMPILER="C:/Qt/Tools/mingw1310_64/bin/g++.exe" ^
    -DCMAKE_MAKE_PROGRAM="C:/Qt/Tools/Ninja/ninja.exe"
ninja
\end{lstlisting}
Adjust the version numbers in the paths to match what you installed
(e.g.\ \texttt{6.10.3}, \texttt{mingw1310\_64}).  Use forward slashes in
paths, even on Windows.

\subsection{CMake cannot find Qt}

If CMake reports that it cannot find Qt6, verify that
\texttt{-DCMAKE\_PREFIX\_PATH} points to the correct Qt installation.
Look for a directory like \texttt{C:\textbackslash Qt\textbackslash 6.10.3\textbackslash mingw\_64}
containing \texttt{lib\textbackslash cmake\textbackslash Qt6}.

\subsection{Wrong compiler picked up}

If CMake finds a different compiler (e.g.\ MSVC), make sure you are
specifying \texttt{-DCMAKE\_C\_COMPILER} and
\texttt{-DCMAKE\_CXX\_COMPILER} explicitly pointing at MinGW's
\texttt{gcc.exe} and \texttt{g++.exe}.

\subsection{DLL not found at runtime}

If \texttt{femm\_gui.exe} fails to launch with missing DLL errors,
add the Qt and MinGW binary directories to your \texttt{PATH}:
\begin{lstlisting}
set PATH=C:\Qt\6.10.3\mingw_64\bin;C:\Qt\Tools\mingw1310_64\bin;%PATH%
bin\femm_gui.exe
\end{lstlisting}

\subsection{Building the standalone mesher}

The \texttt{tangle} mesher can be built independently without Qt,
using just MinGW:
\begin{lstlisting}
cd solvers\common
g++ -O3 -std=c++17 -static -o tangle.exe tangle.cpp float256.cpp -lm
\end{lstlisting}

\end{document}
